\documentclass[10pt,a4paper]{article}
\usepackage[T1]{fontenc}
\usepackage{lmodern,amsmath,amssymb,amsthm,mathtools,microtype,geometry,enumitem,hyperref}
\geometry{left=17mm,right=17mm,top=14mm,bottom=14mm}
\hypersetup{hidelinks,pdftitle={Tanaka's Central Limit Theorem - Revised Joint Course Report}}
\setlength{\parindent}{0pt}\setlength{\parskip}{2.7pt}
\setlist[enumerate]{topsep=4pt,itemsep=3pt}
\newtheorem{theorem}{Theorem}\newtheorem{lemma}{Lemma}
\newcommand{\R}{\mathbb R}\newcommand{\E}{\mathbb E}\newcommand{\Var}{\operatorname{Var}}
\newcommand{\Law}{\operatorname{Law}}\newcommand{\N}{\mathcal N}
\newcommand{\W}{W_2}
\begin{document}
\begin{center}
{\Large\bfseries Tanaka's Central Limit Theorem}\par
{\large A Study in Wasserstein Distance and Coupling}\par
\vspace{10pt}
{\normalsize Probability and Statistics Course Project}\par
{\normalsize Instructor: Dr. Mojahedian}\par
\vspace{5pt}
{\bfseries Soheil Bolvardi \quad Mohammadamin Khedri \quad Houman Rajaeian}\par
\vspace{8pt}
{\small Revised mathematical exposition, October 2026}
\end{center}
\vspace{8pt}
\textbf{Scope and attribution.} The original course report was jointly authored by the three students listed above. This revised exposition corrects mathematical statements and clarifies omitted proof steps while retaining the report's original aim: an undergraduate explanation of Tanaka's central limit theorem through couplings and probability metrics. The proof outline below \emph{cites}, rather than reproves, the nontrivial strictness result in Tanaka's inequality.

\section*{1. Introduction}
The classical central limit theorem states that properly centered and scaled sums of independent and identically distributed random variables converge in distribution to a normal law. Its conclusion is qualitative: weak convergence alone does not require second moments to converge. One can sharpen the language of convergence by introducing a distance between probability laws that depends on both their locations and the amount of mass transported.

The second-order Wasserstein distance is particularly natural for the central limit theorem because its transportation cost is quadratic. The resulting convergence is stronger than weak convergence. However, we should distinguish the existence of a \emph{metric} for weak convergence (several are known) from the specific extra moment information conveyed by convergence in $W_2$.

\begin{theorem}[Central limit theorem in $W_2$]
Let $X_1,X_2,\ldots$ be i.i.d. real random variables with $m=\E X_1$ and $0<\sigma^2=\Var(X_1)<\infty$. For $S_n=X_1+\cdots+X_n$,
\[
\W\left(\Law\left(\frac{S_n-nm}{\sigma\sqrt n}\right),\N(0,1)\right)
\longrightarrow 0 \quad(n\to\infty).
\]
\end{theorem}
The theorem itself assumes only a finite, positive second moment. The coupling-based proof outline in Sections 4--5 uses an \emph{additional finite fourth moment} to justify a passage to a limit in $W_2$. This restriction is a feature of the exposition, not an extra hypothesis of Tanaka's theorem.

\section*{2. Probability and measure-theoretic preliminaries}
On a measurable space $(S,\mathcal S)$, a probability measure $P$ satisfies $P(S)=1$, $P(\varnothing)=0$, and countable additivity on disjoint measurable sets. A probability space is a triple $(\Omega,\mathcal A,P)$. A random variable $X:\Omega\to S$ is measurable, and its \emph{law} is the pushforward measure
\[
\mu_X(B)=P\bigl(X\in B\bigr)=P\bigl(X^{-1}(B)\bigr),\qquad B\in\mathcal S.
\]
This lets us study the distribution of $X$ without describing the underlying sample outcomes individually. The expectation of an integrable real random variable is the Lebesgue integral
\[
\E X=\int_\Omega X\,dP,
\]
and variance is $\Var(X)=\E(X-\E X)^2$. Monotone and dominated convergence theorems explain circumstances in which limits of random variables can be passed through integrals or expectations.

We will need the dominated-convergence principle in its correct form: if $Y_n\to Y$ almost surely and $|Y_n|\le Z$ almost surely for one integrable $Z$, then $\E Y_n\to\E Y$. More generally, the fourth-moment bound used in the coupling argument implies \emph{uniform integrability} of the squared normalized sums. That fact supplies convergence of their second moments along a weakly convergent subsequence.

\section*{3. Distances between distributions and couplings}
The \emph{total variation distance} of probability measures on the same measurable space is
\[
d_{\mathrm{TV}}(\mu,\nu)=\sup_{A\in\mathcal S}|\mu(A)-\nu(A)|.
\]
When $\mu$ and $\nu$ have densities $f,g$ with respect to the same reference measure $\lambda$, it equals $\frac12\int|f-g|\,d\lambda$. It measures the greatest possible difference in event probabilities.

A \emph{coupling} of laws $\mu$ and $\nu$ is a pair $(U,V)$ defined on a common probability space with the specified marginal laws. It may introduce dependence between the components without changing either marginal. Total variation admits the coupling characterization
\[
d_{\mathrm{TV}}(\mu,\nu)=\inf\bigl\{P(U\ne V):(U,V)\text{ couples }\mu,\nu\bigr\}
\]
(on standard Borel spaces), motivating the idea of comparing distributions by optimizing over joint realizations.

For $p\ge1$ and laws with finite $p$th moments, the \emph{Wasserstein distance} is
\[
W_p(\mu,\nu)
=\left(\inf_{(U,V)\in\Gamma(\mu,\nu)}\E|U-V|^p\right)^{1/p}.
\]
For probability laws on the real line, the generalized inverse distribution functions $Q_\mu,Q_\nu$ give the quantile-coupling formula
\[
W_p^p(\mu,\nu)=\int_0^1|Q_\mu(u)-Q_\nu(u)|^p\,du.
\]
The absolute value is necessary for a correct distance formula. It is this optimal-coupling representation, specialized to $p=2$, that underlies the quadratic comparison inequality below.


\subsection*{An illustrative coupling comparison}
The difference between total variation and Wasserstein costs can be seen already in discrete examples. If $X\sim\mathrm{Ber}(p)$ and $Y\sim\mathrm{Ber}(q)$, any coupling satisfies
\[
d_{\mathrm{TV}}(\Law(X),\Law(Y))\le P(X\ne Y).
\]
The optimal coupling achieves equality, namely $|p-q|$. For the independent coupling, by contrast,
\[
P(X\ne Y)=p(1-q)+(1-p)q=p+q-2pq,
\]
which can be much larger. This example illustrates why the comparison of distributions requires optimizing over joint laws rather than assuming the observations are independently paired.

\section*{4. Tanaka's comparison inequality}
Suppose independent, centered random variables $X,Y$ have finite variances $\sigma_X^2,\sigma_Y^2$. Let $Z_X,Z_Y$ be independent centered Gaussians with those variances, optimally coupled to $X$ and $Y$ in two independent pairings. For arbitrary real coefficients $a,b$, let
\[
Z=aZ_X+bZ_Y\sim\N\left(0,a^2\sigma_X^2+b^2\sigma_Y^2\right).
\]
Because the coupling errors have zero means and are independent,
\begin{align*}
&\W^2\bigl(\Law(aX+bY),\Law(Z)\bigr)\\
&\quad\le\E|a(X-Z_X)+b(Y-Z_Y)|^2\\
&\quad=a^2\W^2\bigl(\Law(X),\Law(Z_X)\bigr)
+b^2\W^2\bigl(\Law(Y),\Law(Z_Y)\bigr).
\end{align*}
This is the basic improvement bound: combining independent centered components does not increase the Gaussian comparison cost beyond the corresponding weighted sum. The stronger \emph{strictness/equality characterization} for nonzero coefficients, used next, is a separate nontrivial result in Tanaka's work. We do not derive it from linearity of normal quantile functions: normal quantile functions are nonlinear.

\section*{5. Coupling-based proof outline under a finite fourth moment}
Assume for this section that $\E X_1=0$, $\E X_1^2=1$, and $\E X_1^4<\infty$, and set
\[
G_n=\frac{1}{\sqrt n}\sum_{i=1}^nX_i,\qquad
\eta_k=G_{2^k},\qquad
D_k=\W^2\bigl(\Law(\eta_k),\N(0,1)\bigr).
\]
Let $\eta_k'$ be an independent copy of $\eta_k$. Since
$\eta_{k+1}\stackrel d=(\eta_k+\eta_k')/\sqrt2$, the comparison inequality yields
\[
0\le D_{k+1}\le D_k.
\]
Therefore $D_k$ converges to some $\ell\ge0$. The fourth-moment assumption gives
\[
\E G_n^4=\frac{\E X_1^4}{n}+3\left(1-\frac1n\right),
\]
so the fourth moments are uniformly bounded. Tightness (from variance one) gives a weakly convergent subsequence of $\Law(\eta_k)$; uniform integrability of the squares upgrades that subsequence to convergence in $W_2$ to some law $\mu$ with mean zero and variance one.

Let $U,V$ be independent with law $\mu$. By continuity in $W_2$ along the subsequence and its one-step successors,
\[
\W^2(\mu,\N(0,1))=\ell
=\W^2\left(\Law\left(\frac{U+V}{\sqrt2}\right),\N(0,1)\right).
\]
The cited strictness/equality characterization implies $\mu=\N(0,1)$; hence $\ell=0$. Thus convergence in $W_2$ is established on the dyadic subsequence $n=2^k$.

For an arbitrary $n$, write its binary expansion $n=\sum_{k\ge0}\varepsilon_k2^k$ with $\varepsilon_k\in\{0,1\}$, grouping the original $X_i$ into independent blocks. Applying the comparison inequality gives
\[
\W^2(\Law(G_n),\N(0,1))
\le\sum_{k\ge0}\varepsilon_k\frac{2^k}{n}D_k.
\]
For any $\epsilon>0$, choose $K$ with $D_k<\epsilon$ for all $k\ge K$. As $D_k\le D_0$, the right side is bounded by $D_0(2^K-1)/n+\epsilon$, proving convergence to zero as $n\to\infty$. This argument \emph{does not} show monotonicity of the distance for every integer $n$; only the dyadic subsequence was shown to be nonincreasing.

Finally, the full theorem with only $\E X_1^2<\infty$ can also be obtained from the classical central limit theorem and the standard criterion that weak convergence together with convergence of second moments is equivalent to $W_2$ convergence on $\R$. Each normalized sum has second moment one, equal to that of the normal limit. This is an alternative deduction of the full conclusion, not a self-contained proof of Tanaka's strictness result.


\subsection*{What is proved and what is quoted}
The coupling argument has several distinct components. First, optimal pairings with matching normal laws yield a quadratic comparison bound. Second, the squared Wasserstein distance of dyadic normalized sums decreases. Third, uniform fourth-moment control allows a subsequential weak limit to be upgraded to a limit in $W_2$. Finally, Tanaka's strictness result, \emph{used as an external theorem}, forces any such limiting law to be Gaussian. The binary-block estimate then extends convergence beyond powers of two.

It would be incorrect to claim that the quadratic comparison inequality alone proves the strictness characterization. Nor does monotonicity of dyadic distances mean that the distance decreases at \emph{every} successive sample size. These distinctions are important to the correct scope of the reading project and the cited research.

\section*{6. Conclusion}
Couplings and probability metrics make it possible to compare distributions quantitatively. In the coupling outline, quadratic transport costs decrease along dyadic normalized sums, and the strictness of Tanaka's inequality identifies possible limiting distributions as Gaussian. Binary decomposition extends the dyadic result to all sample sizes. The measure-theoretic preliminaries clarify the additional hypotheses needed to justify limit passages. The theorem is stronger than weak convergence, although it does not by itself provide a numerical convergence \emph{rate}.

\section*{References}
\begin{enumerate}[leftmargin=*]
\item H. Tanaka (1973), ``An inequality for a functional of probability distributions and its application to Kac's one-dimensional model of a Maxwellian gas,'' \emph{Zeitschrift f\"ur Wahrscheinlichkeitstheorie und Verwandte Gebiete}, 27, 47--52.
\item O. Johnson and R. Samworth (2005), ``Central limit theorem and convergence to stable laws in Mallows distance,'' \emph{Bernoulli}, 11(5), 829--845.
\item L. Leskel\"a, \emph{Probability Theory: A Fast Course}, especially \S12.5.
\end{enumerate}
\end{document}
